\chapter{LITERATURE REVIEW}\label{chap-lit}
\thispagestyle{empty}

In this chapter, we review the existing literature related to the design and implementation of cold plasma generators, high-voltage pulsed power supplies, and plasma jet applications. The aim is to identify existing solutions, evaluate their performance, and find gaps that justify the development of the present project.\\
The literature review is organized by classifying the available knowledge into sections and sub-sections, highlighting different design approaches, and pointing out limitations and opportunities for improvement.

\section{High-Voltage Cold Plasma Generators}\label{sec-rev-topic1}

Cold plasma generators rely on high-voltage, high-frequency pulsed power supplies to ionize a carrier gas such as argon or helium. The main challenge is achieving stable plasma generation at atmospheric pressure without excessive heating.

\subsection{Laboratory-Scale Designs}\label{sec-rev-topic11}

Diaz et al.~\cite{Diaz2022} presented a low-cost cold plasma generator designed for laboratory applications. The design employed an H-bridge inverter driving a high-voltage transformer, followed by a Cockcroft--Walton voltage multiplier to generate DC voltages in the 20--30~kV range. While the system was effective for small-scale experiments, it lacked precise voltage control and advanced switching features.

\subsection{Compact High-Voltage Pulse Generators}\label{sec-rev-topic12}

Spassov et al.~\cite{Spassov1997} demonstrated a compact high-voltage pulse generator for plasma applications. The system achieved reliable plasma generation, but its switching frequency was limited and the design was not optimized for atmospheric pressure cold plasma jets.

\subsection{Atmospheric Pressure Cold Plasma Jets}\label{sec-rev-topic13}

Cheng et al.~\cite{Cheng2006} developed an atmospheric pressure cold plasma jet generator for sterilization applications. The generator utilized high-voltage pulses to create a visible plasma jet. While the device successfully sterilized surfaces, the design required complex circuitry and offered limited output power control.




\section{Methodology}\label{sec-rm}

The methodology adopted for this project includes the following steps:  
\begin{enumerate}
    \item Studying literature on cold plasma generators, H-bridge inverters, high-voltage transformers, and Cockcroft--Walton voltage multipliers.
    \item Designing the H-bridge inverter and selecting suitable MOSFETs (IRFP460) and gate drivers (UCC21330).
    \item Programming the STM32G431 Microcontroller to generate high-resolution PWM signals for precise control of the H-bridge.
    \item Constructing a prototype and integrating it with a mechanical setup for argon plasma jet generation.
    \item Measuring output voltage, frequency, and plasma characteristics to validate design objectives.
\end{enumerate}

\section{Summary}

This chapter reviewed the existing designs of cold plasma generators and high-voltage pulse systems. The analysis identified limitations in voltage control, switching performance, and integration with plasma jets. These findings form the basis for designing a ``compact, controllable cold plasma generator'' using the STM32G431 MCU, IRFP460 MOSFET H-bridge, and Cockcroft--Walton voltage multiplier for laboratory applications.

% References cited in this chapter
\nocite{Diaz2022, Spassov1997, Cheng2006}
